% =====================================================================
\chapter{Soluciones de los ejercicios}
\label{ap:soluciones}
% =====================================================================

Cada solución se presenta en cinco partes: datos, planteamiento, desarrollo, resultado y comprobación. En los pasos intermedios se redondea a cuatro cifras significativas salvo que se indique otra cosa. Los resultados numéricos de los casos de estudio se han obtenido y comprobado con el cuaderno \texttt{M28-practica.ipynb}.

% ---------------------------------------------------------------------
\section{Tema 1}

\begin{solucion}{eje:1-componentes}
\datos cola de una máquina; trabajos con tiempo de proceso $p_j$, fecha de entrega $d_j$ y peso $w_j$.
\planteamiento identificar los elementos de la \cref{def:busqueda-programas} y la \cref{def:esqueleto}.
\pasos
\begin{enumerate}[label=\arabic*.]
  \item Esqueleto: un simulador de la máquina que, cada vez que queda libre, calcula la prioridad de cada trabajo en cola y procesa el de mayor prioridad; garantiza que cada trabajo se procesa una vez y sin solaparse.
  \item Firma: \texttt{prioridad(p, d, w, t\_actual) -> float}, vectorizada sobre los trabajos en cola.
  \item Utilidad: $u(\pi,x)=-\sum_j w_j\max(0,C_j-d_j)$, el retraso ponderado cambiado de signo, con $C_j$ el instante de fin; $-\infty$ si el programa falla.
  \item Entrenamiento: varias semanas de llegadas históricas (o simuladas con su distribución); prueba: semanas posteriores y, aparte, semanas con más carga.
  \item Semillas: la regla del menor tiempo de proceso ponderado, $w_j/p_j$, y la de la fecha de entrega más temprana, $-d_j$.
\end{enumerate}
\resultado \textbf{esqueleto = simulador de la cola; función que evoluciona = prioridad; $u$ = menos el retraso ponderado; semillas $w_j/p_j$ y $-d_j$.}
\comprobacion cualquier función numérica produce una secuencia factible, como exige la \cref{def:esqueleto}.
\end{solucion}

\begin{solucion}{eje:1-softmax}
\datos $s=(-416{,}7;\ -409{,}3;\ -405{,}3)$; $T=0{,}1$ sin normalizar; $T=0{,}3$ con normalización.
\planteamiento \cref{def:softmax-clusters}; restar el máximo para evitar desbordamientos.
\pasos
\begin{enumerate}[label=\arabic*.]
  \item (a) $(s-\max s)/T=(-114;\ -40;\ 0)$; $e^{-114}\approx 3\cdot10^{-50}$ y $e^{-40}\approx4\cdot10^{-18}$, luego $p\approx(0;\ 0;\ 1)$.
  \item (b) Rango $11{,}4$: $z=(-1;\ -0{,}3509;\ 0)$; $z/T=(-3{,}333;\ -1{,}170;\ 0)$; exponenciales $(0{,}0357;\ 0{,}3105;\ 1)$, suma $1{,}3462$; $p=(0{,}0265;\ 0{,}2307;\ 0{,}7428)$.
  \item (c) La temperatura solo tiene sentido relativa a la escala de las puntuaciones: con puntuaciones del orden de cientos, $T=0{,}1$ hace el muestreo voraz; normalizando, la misma $T$ significa lo mismo en cualquier problema.
\end{enumerate}
\resultado \textbf{(a) $(0;0;1)$ en la práctica; (b) $(0{,}027;\ 0{,}231;\ 0{,}743)$; (c) la escala debe fijarse antes de elegir $T$.}
\comprobacion $0{,}0265+0{,}2307+0{,}7428=1$; por la \cref{prop:temperatura}, al bajar $T$ la probabilidad se concentra en el mejor.
\end{solucion}

\begin{solucion}{eje:1-ffbf}
\datos $C=10$; objetos $3,7,4,7,8,3,5$; $S=37$.
\planteamiento aplicar la \cref{def:ff-bf} y el \cref{lem:l1}.
\pasos
\begin{enumerate}[label=\arabic*.]
  \item $L_1=\lceil 3{,}7\rceil=4$.
  \item FF: $3\to$1; $7\to$1 (carga 10); $4\to$2; $7\to$3; $8\to$4; $3\to$2 (carga 7); $5$: no cabe en 2 (resto 3), 3 (resto 3) ni 4 (resto 2) y abre el 5. Total 5.
  \item BF: igual hasta el segundo $3$, que cabe en 2 (hueco 3), 3 (hueco 0) y 4 (no cabe, resto 2): va al 3 (carga 10); el $5$ cabe en 2 (resto 6, hueco 1) y va allí. Total 4.
\end{enumerate}
\resultado \textbf{FF usa 5 contenedores y BF 4; BF es óptimo porque alcanza $L_1=4$.}
\comprobacion cargas de BF: $\{3,7\}$, $\{4,5\}$, $\{7,3\}$, $\{8\}$: suman $10+9+10+8=37$.
\end{solucion}

\begin{solucion}{eje:1-nextfit}
\datos NF con capacidad $C$; ejemplo con $C=1$ y objetos $1/2,\varepsilon,1/2,\varepsilon,\dots$ ($2k$ objetos).
\planteamiento cotas por pares de contenedores consecutivos y \cref{lem:l1}.
\pasos
\begin{enumerate}[label=\arabic*.]
  \item (a) Si NF cerró el contenedor $i$ para abrir el $i+1$, el primer objeto de $i+1$ no cabía en $i$: carga$(i)$ $+$ carga$(i+1)>C$.
  \item Con $m$ contenedores, los $\lfloor m/2\rfloor$ pares disjuntos $(1,2),(3,4),\dots$ suman más de $C$ cada uno: $S>\lfloor m/2\rfloor C$. Como $S\le\mathrm{OPT}\cdot C$, se tiene $\lfloor m/2\rfloor<\mathrm{OPT}$ y, por ser enteros, $\lfloor m/2\rfloor\le\mathrm{OPT}-1$. Por tanto $m\le2\lfloor m/2\rfloor+1\le2\,\mathrm{OPT}-1$.
  \item (b) NF: el contenedor con $1/2$ admite $\varepsilon$ (carga $1/2+\varepsilon$), pero el siguiente $1/2$ ya no cabe; cada par $(1/2,\varepsilon)$ ocupa un contenedor: $k$ contenedores.
  \item Óptimo: los $k$ objetos de $1/2$ caben de dos en dos en $\lceil k/2\rceil$ contenedores, y los $k$ objetos $\varepsilon$ (total $k\varepsilon<1/2$) en uno más (o en el hueco libre si $k$ es impar): $\mathrm{OPT}\le k/2+1$.
  \item Cociente: $k/(k/2+1)\to 2$ cuando $k\to\infty$.
\end{enumerate}
\resultado \textbf{NF $\le 2\,\mathrm{OPT}-1$, y el factor 2 es asintóticamente alcanzable: $k$ frente a $k/2+1$.}
\comprobacion con $k=10$ y $\varepsilon=0{,}01$: NF usa 10 contenedores; el óptimo, 6 (cinco con dos mitades y uno con los diez $\varepsilon$, carga $0{,}1$); $10\le 2\cdot6$.
\end{solucion}

\begin{solucion}{eje:1-presupuesto}
\datos 2 mejoras en 1000 propuestas; 0{,}004~€ por llamada; 20~s de CPU a 0{,}05~€/h; 10 islas de 300 evaluaciones.
\planteamiento \cref{prop:espera} (ii) y (iii).
\pasos
\begin{enumerate}[label=\arabic*.]
  \item (a) $\hat q=2/1000=0{,}002$.
  \item (b) $(1-q)^{B}\le0{,}05\iff B\ge\ln0{,}05/\ln0{,}998=1496{,}4$: $B=1497$. Con la cota $e^{-qB}\le0{,}05$ se obtiene $B\ge\ln20/0{,}002=1497{,}9$.
  \item (c) LLM: $1497\cdot0{,}004=5{,}99$~€; CPU: $1497\cdot20/3600\cdot0{,}05=0{,}42$~€; total $6{,}40$~€.
  \item (d) Cada isla falla con probabilidad $0{,}998^{300}=0{,}5485$; las diez, $0{,}5485^{10}=0{,}00246$; alguna mejora con probabilidad $0{,}9975$.
\end{enumerate}
\resultado \textbf{$q\approx0{,}002$; $B=1497$ evaluaciones; unos $6{,}40$~€; $0{,}9975$ con 10 islas.}
\comprobacion (d) coincide con $1-0{,}998^{3000}$: con independencia, repartir las evaluaciones en islas no cambia la probabilidad de alguna mejora; las islas ayudan por la diversidad, que este modelo no recoge.
\end{solucion}

\begin{solucion}{eje:1-gramatica}
\datos $E\to E+E\mid E\times E\mid x\mid 1$.
\planteamiento contar árboles por la raíz.
\pasos
\begin{enumerate}[label=\arabic*.]
  \item (a) Un árbol de profundidad $\le d$ es una hoja (2 opciones) o una raíz $+$ o $\times$ (2 opciones) con dos subárboles de profundidad $\le d-1$ ($N_{d-1}$ opciones cada uno): $N_d=2+2N_{d-1}^{2}$, con $N_0=2$.
  \item (b) $N_1=2+2\cdot4=10$; $N_2=2+2\cdot100=202$; $N_3=2+2\cdot202^{2}=81\,610$.
  \item (c) Si solo uno de los $81\,610$ árboles mejora la semilla, $q=1/81\,610=1{,}23\cdot10^{-5}$ y $\E[N]\approx81\,610$ evaluaciones.
\end{enumerate}
\resultado \textbf{$N_1=10$, $N_2=202$, $N_3=81\,610$; $q\approx1{,}2\cdot10^{-5}$.}
\comprobacion el crecimiento es doblemente exponencial ($N_d\approx 2N_{d-1}^{2}$); por eso una distribución de propuestas informada (un LLM) importa tanto.
\end{solucion}

\begin{solucion}{eje:1-interpretar}
\datos $\pi^{*}=-\min(g,9)-5{,}82\,\Ind[g\le20]$.
\planteamiento estudiar la función por intervalos de $g$ y aplicar la \cref{def:ff-bf}.
\pasos
\begin{enumerate}[label=\arabic*.]
  \item (a) $0\le g\le9$: $-g-5{,}82\in[-14{,}82;-5{,}82]$; $9<g\le20$: $-14{,}82$; $g>20$: $-9$. Orden: $g\le3$ (mayor que $-9$; se prefiere el menor $g$), después $g>20$ (empate en $-9$), después $3<g\le9$ (de $-8{,}82$ a $-14{,}82$, decreciente; para $g\in(3{,}18;\,9]$ por debajo de $-9$) y por último $9<g\le20$. Precisando: $-g-5{,}82>-9\iff g<3{,}18$.
  \item (b) Si ningún contenedor deja $g\le20$, todos puntúan $-9$ y el esqueleto elige el primero (desempate por índice): FF. Si alguno deja $g\le3$, ese grupo tiene las mayores puntuaciones y, dentro de él, gana el de menor $g$: BF restringido a esos contenedores.
  \item (c) El contenedor vacío tiene $g=C-s$. Si el objeto cabe en un contenedor abierto con $g\le 20$ o $g>20$, su puntuación es al menos $-14{,}82$, y el vacío puntúa $-9$ si $C-s>20$. Luego $\pi^{*}$ \emph{puede} abrir un contenedor nuevo con un abierto disponible si todos los abiertos dejan $3{,}18<g\le20$ y $s<80$: \emph{no} es \emph{any fit} en general. Por tanto el \cref{teo:anyfit} no se aplica directamente; solo valdría si el esqueleto obligara a usar un abierto cuando el objeto cabe.
\end{enumerate}
\resultado \textbf{preferencia: $g<3{,}18$, luego $g>20$ (FF), luego $3{,}18<g\le9$, luego $9<g\le20$; coincide con FF o BF en los casos descritos; no es \emph{any fit}, de modo que no hereda la garantía.}
\comprobacion con $s=50$ y un único abierto con $r=60$ ($g=10$, puntuación $-14{,}82$), el vacío ($g=50$) puntúa $-9$ y gana: se abre un contenedor teniendo hueco. La heurística paga ese contenedor para no dejar un resto de 10 unidades.
\end{solucion}

\begin{solucion}{eje:1-diseno}
\datos colocación de máquinas virtuales con CPU y memoria.
\planteamiento recorrer los componentes de FunSearch y los errores de la \cref{obs:errores1}.
\pasos
\begin{enumerate}[label=\arabic*.]
  \item Esqueleto: simulador de colocación en línea; para cada máquina virtual, calcula los servidores factibles (CPU y memoria libres suficientes, reglas de antiafinidad) y llama a \texttt{prioridad(cpu, mem, cpu\_libre, mem\_libre)}.
  \item Evaluador: recursos varados medios (capacidad libre de un recurso en servidores con el otro agotado), número de servidores encendidos y peticiones rechazadas; puntuación escalar con pesos acordados; ejecución en contenedor sin red, con tiempo y memoria limitados; el programa no recibe las trazas completas, solo los argumentos de la firma.
  \item Datos: trazas por semanas y clústeres; entrenamiento con unas semanas, validación con otras para elegir el programa y prueba con semanas y clústeres no usados, más escenarios de carga alta.
  \item Presupuesto: piloto de 1000 propuestas para estimar $q$; $B\approx 3/q$ para un 95\,\% de probabilidad de mejora (\cref{prop:espera}); coste estimado antes de lanzar.
  \item Despliegue: simplificar y leer el programa; comparar con BF por recurso dominante con intervalos; despliegue en sombra y luego gradual con la regla actual como respaldo.
\end{enumerate}
\resultado \textbf{esqueleto factible por construcción, evaluador aislado y multiobjetivo, división por tiempo y clúster, presupuesto con $1/q$ y despliegue gradual.}
\comprobacion cada riesgo de la \cref{obs:errores1} tiene una medida asociada.
\end{solucion}

% ---------------------------------------------------------------------
\section{Tema 2}

\begin{solucion}{eje:2-meta}
\datos clasificación de partes de avería en seis categorías.
\planteamiento \cref{def:opro}.
\pasos
\begin{enumerate}[label=\arabic*.]
  \item Cabecera: descripción de la tarea (clasificar un parte en una de seis categorías) y tres ejemplos de partes con su categoría.
  \item Trayectoria: las 20 mejores instrucciones con su exactitud de validación, \emph{de menor a mayor}, de modo que las mejores queden al final, junto a la petición.
  \item Petición: «escribe una instrucción nueva, distinta de las anteriores, que obtenga mayor exactitud».
  \item No debe contener: ejemplos del conjunto de prueba, la respuesta correcta de los ejemplos de validación que se usan para puntuar, ni datos personales de los partes.
\end{enumerate}
\resultado \textbf{tarea y ejemplos, trayectoria ordenada de peor a mejor, petición; sin datos de prueba ni datos personales.}
\comprobacion el orden ascendente sigue la descripción del método~\cite{opro}.
\end{solucion}

\begin{solucion}{eje:2-metrica}
\datos sistema de preguntas y respuestas con recuperación de fragmentos.
\planteamiento \cref{def:dspy}: la métrica guía qué trazas se conservan.
\pasos
\begin{enumerate}[label=\arabic*.]
  \item Métrica propuesta:
  \[
  m(e,\hat e)=\Ind[\text{respuesta correcta}]\cdot\Ind[\text{algún fragmento recuperado contiene la respuesta}],
  \]
  con la corrección medida por coincidencia normalizada o por $F_1$ de palabras.
  \item Con solo la respuesta, el optimizador conservaría como demostraciones trazas que aciertan por casualidad o por conocimiento previo del modelo, sin usar el contexto; el programa aprendería a ignorar la recuperación y fallaría en manuales que el modelo no conoce.
\end{enumerate}
\resultado \textbf{exigir a la vez respuesta correcta y contexto que la respalde.}
\comprobacion es la idea de la función de validación de contexto y respuesta de la tarea de CS224U~\cite{cs224u}.
\end{solucion}

\begin{solucion}{eje:2-maldicion}
\datos $K=50$; $n=200$; $J=0{,}80$.
\planteamiento \cref{prop:maldicion} (ii) con $\sigma=\sqrt{J(1-J)/n}$.
\pasos
\begin{enumerate}[label=\arabic*.]
  \item (a) $\sigma=\sqrt{0{,}16/200}=0{,}02828$; $\sqrt{2\ln50}=2{,}797$; cota $0{,}0791$.
  \item (b) $\sqrt{0{,}16/n}\cdot2{,}797\le0{,}02\iff n\ge0{,}16\cdot2\ln50/0{,}0004=3130$.
  \item (c) Con $K=5$: $\sqrt{2\ln5}=1{,}794$; cota con $n=200$: $0{,}0507$; para 2 puntos, $n\ge0{,}16\cdot2\ln5/0{,}0004=1288$.
\end{enumerate}
\resultado \textbf{(a) $\sigma=0{,}0283$, cota $0{,}079$; (b) 3130 ejemplos; (c) cota $0{,}051$ y 1288 ejemplos.}
\comprobacion la cota crece como $\sqrt{\ln K}$: pasar de 5 a 50 candidatos multiplica el tamaño necesario por $\ln50/\ln5=2{,}43$, y $3130/1288=2{,}43$.
\end{solucion}

\begin{solucion}{eje:2-voto}
\datos $p=0{,}7$; $V=0{,}20$~€; $c=0{,}004$~€.
\planteamiento \cref{prop:incremento} y \cref{cor:m-optimo}: seguir mientras $V\Delta_m\ge2c=0{,}008$.
\pasos
\begin{enumerate}[label=\arabic*.]
  \item $p(1-p)=0{,}21$, $2p-1=0{,}4$. $\Delta_1=\binom10\cdot0{,}21\cdot0{,}4=0{,}0840$; $\Delta_3=\binom31\cdot0{,}21^{2}\cdot0{,}4=0{,}05292$; $\Delta_5=\binom52\cdot0{,}21^{3}\cdot0{,}4=0{,}03704$; $\Delta_7=\binom73\cdot0{,}21^{4}\cdot0{,}4=0{,}02723$.
  \item $V\Delta_1=0{,}0168\ge0{,}008$; $V\Delta_3=0{,}0106\ge0{,}008$; $V\Delta_5=0{,}00741<0{,}008$: $m^{*}=5$.
  \item $P_5=P_1+\Delta_1+\Delta_3=0{,}7+0{,}084+0{,}0529=0{,}8369$.
\end{enumerate}
\resultado \textbf{$m^{*}=5$, con $P_5=0{,}837$ y beneficio $0{,}2\cdot0{,}8369-0{,}02=0{,}1474$~€.}
\comprobacion beneficios con $m=3$ y $m=7$: $0{,}2\cdot0{,}784-0{,}012=0{,}1448$ y $0{,}2\cdot0{,}8740-0{,}028=0{,}1468$, ambos menores.
\end{solucion}

\begin{solucion}{eje:2-amdahl}
\datos $f=0{,}7$ (bucle), $0{,}2$ (lectura), $0{,}1$ (resto).
\planteamiento \cref{prop:amdahl}.
\pasos
\begin{enumerate}[label=\arabic*.]
  \item (a) $A=1/(0{,}3+0{,}7/8)=1/0{,}3875=2{,}58$.
  \item (b) $1/(0{,}3+0{,}7/a)=2\iff0{,}7/a=0{,}2\iff a=3{,}5$.
  \item (c) Con $a\to\infty$, $A\to1/0{,}3=3{,}33<4$: imposible sin acelerar también la lectura o el resto.
\end{enumerate}
\resultado \textbf{(a) $2{,}58$; (b) factor $3{,}5$; (c) no: el máximo es $3{,}33$.}
\comprobacion con $a=3{,}5$: $0{,}3+0{,}2=0{,}5$ y $A=2$.
\end{solucion}

\begin{solucion}{eje:2-medida}
\datos $t_A=(2;4;1;8;3)$, $t_B=(1{,}6;4{,}4;0{,}5;8;2{,}4)$; 14 victorias de B en 20 pares.
\planteamiento medias aritmética y geométrica de cocientes; contraste de signos.
\pasos
\begin{enumerate}[label=\arabic*.]
  \item (a) $t_B/t_A=(0{,}8;\ 1{,}1;\ 0{,}5;\ 1;\ 0{,}8)$: media aritmética $0{,}840$, geométrica $(0{,}8\cdot1{,}1\cdot0{,}5\cdot1\cdot0{,}8)^{1/5}=0{,}352^{1/5}=0{,}8115$.
  \item (b) $t_A/t_B=(1{,}25;\ 0{,}909;\ 2;\ 1;\ 1{,}25)$: aritmética $1{,}282$, geométrica $1{,}2322$.
  \item (c) Las aritméticas dicen «B tarda un 84\,\% de A» y «A tarda un 128\,\% de B», lo que implicaría $0{,}84\cdot1{,}28=1{,}08\ne1$: incoherente. Las geométricas cumplen $0{,}8115\cdot1{,}2322=1$, porque la media geométrica de los inversos es el inverso de la media geométrica.
  \item (d) $\Prob(\text{Bin}(20;0{,}5)\ge14)=0{,}0577>0{,}05$: no se alcanza el criterio de 15 victorias.
\end{enumerate}
\resultado \textbf{media geométrica: B tarda un 81{,}2\,\% de A; con 14 de 20 victorias no se acepta.}
\comprobacion $0{,}8115^{-1}=1{,}2323$.
\end{solucion}

\begin{solucion}{eje:2-cota}
\datos $\varepsilon_k$ independientes $\mathcal{N}(0,\sigma^{2})$.
\planteamiento monotonía del máximo; para $K=2$, la identidad $\max(a,b)=\frac{a+b}2+\frac{\abs{a-b}}2$.
\pasos
\begin{enumerate}[label=\arabic*.]
  \item $\max_{k\le K+1}\varepsilon_k\ge\max_{k\le K}\varepsilon_k$ para cada realización; tomando esperanzas, $\E[\max]$ es no decreciente en $K$.
  \item $K=2$: $\E[(\varepsilon_1+\varepsilon_2)/2]=0$ y $\varepsilon_1-\varepsilon_2\sim\mathcal{N}(0,2\sigma^{2})$; para $Z\sim\mathcal{N}(0,\tau^{2})$, $\E\abs{Z}=\tau\sqrt{2/\pi}$. Con $\tau=\sigma\sqrt2$: $\E\abs{\varepsilon_1-\varepsilon_2}=2\sigma/\sqrt\pi$.
  \item $\E[\max]=\frac12\cdot\frac{2\sigma}{\sqrt\pi}=\sigma/\sqrt\pi=0{,}5642\,\sigma$.
  \item Cota: $\sigma\sqrt{2\ln2}=1{,}1774\,\sigma$.
\end{enumerate}
\resultado \textbf{$\E[\max(\varepsilon_1,\varepsilon_2)]=\sigma/\sqrt\pi\approx0{,}564\,\sigma$, frente a la cota $1{,}177\,\sigma$.}
\comprobacion el valor exacto es positivo y menor que la cota, como exige la \cref{prop:maldicion}.
\end{solucion}

\begin{solucion}{eje:2-diseno}
\datos asistente que redacta órdenes de trabajo; optimización con DSPy.
\planteamiento aplicar la \cref{def:dspy}, la \cref{prop:maldicion} y el \cref{cor:protocolo}.
\pasos
\begin{enumerate}[label=\arabic*.]
  \item Datos: partes históricos con su orden de trabajo aprobada; entrenamiento (para demostraciones), validación (para elegir) y prueba (reservada), divididos por fecha.
  \item Métrica: campos obligatorios correctos (equipo, tipo de intervención, repuestos) con comprobación automática, más una valoración de redacción por un LLM juez; el juez se valida contra 100 valoraciones de técnicos (acuerdo y sesgos) antes de usarlo.
  \item Presupuesto: número máximo de llamadas por compilación y coste estimado; se detiene si dos rondas no mejoran la validación más allá de su error típico.
  \item Aceptación: en el conjunto de prueba, intervalo de confianza de la diferencia pareada entre la versión nueva y la actual sobre los mismos partes; se acepta si el límite inferior es positivo y ningún campo obligatorio empeora.
  \item Cambio de modelo del proveedor: fijar la versión del modelo; ante una actualización, repetir la evaluación de prueba y recompilar si empeora.
\end{enumerate}
\resultado \textbf{división temporal en tres partes, métrica verificable con juez validado, presupuesto con parada, aceptación por intervalo pareado y reevaluación ante cambios de modelo.}
\comprobacion el protocolo evita los errores (a)--(d) de la \cref{obs:errores2}.
\end{solucion}

% ---------------------------------------------------------------------
\section{Tema 3}

\begin{solucion}{eje:3-dimensiones}
\datos cuatro sistemas.
\planteamiento las tres dimensiones de Elsken \emph{et al.}~\cite{elsken}.
\pasos
\begin{enumerate}[label=\arabic*.]
  \item (a) ENAS: espacio de celdas (o de redes recurrentes) como subgrafos de un supergrafo; estrategia: controlador entrenado con refuerzo; estimación: pesos compartidos.
  \item (b) DARTS: espacio de celdas relajado de forma continua; estrategia: descenso de gradiente binivel; estimación: el propio supergrafo entrenado, con pesos compartidos.
  \item (c) Real \emph{et al.}: espacio de celdas; estrategia: evolución por torneo con envejecimiento; estimación: entrenamiento de cada candidata con presupuesto reducido.
  \item (d) Caso: espacio CASH de tres familias; estrategia: muestreo aleatorio; estimación: multifidelidad por fracción de datos.
\end{enumerate}
\resultado \textbf{ver la lista.}
\comprobacion las tres dimensiones son independientes: cada sistema combina una opción de cada una.
\end{solucion}

\begin{solucion}{eje:3-rejilla}
\datos 6 hiperparámetros, 2 importantes; 64 evaluaciones.
\planteamiento \cref{obs:aleatoria}.
\pasos
\begin{enumerate}[label=\arabic*.]
  \item (a) Rejilla $2^{6}=64$: cada hiperparámetro toma 2 valores; los dos importantes forman $2\cdot2=4$ combinaciones, repetidas 16 veces cada una.
  \item (b) Aleatoria: con probabilidad 1, 64 valores distintos de cada uno y 64 combinaciones distintas de los dos importantes.
  \item (c) $1-0{,}95^{n}\ge0{,}95\iff n\ge\ln0{,}05/\ln0{,}95=58{,}4$: $n=59$.
\end{enumerate}
\resultado \textbf{(a) 2 valores y 4 combinaciones; (b) 64 y 64; (c) 59 puntos.}
\comprobacion $1-0{,}95^{59}=0{,}9515$ y $1-0{,}95^{58}=0{,}9490$.
\end{solucion}

\begin{solucion}{eje:3-cash}
\datos tres algoritmos con hiperparámetros condicionales.
\planteamiento \cref{def:cash} con un hiperparámetro raíz.
\pasos
\begin{enumerate}[label=\arabic*.]
  \item Raíz: algoritmo $\in\{k\text{NN},\ \text{SVM},\ \text{bosque}\}$.
  \item Rama $k$NN: $k\in\{1,\dots,50\}$, distancia $\in\{\text{euclídea},\ \text{Manhattan}\}$.
  \item Rama SVM: núcleo $\in\{\text{lineal},\ \text{gaussiano}\}$, $C>0$ (escala logarítmica); si núcleo $=$ gaussiano, $\gamma>0$ (escala logarítmica).
  \item Rama bosque: número de árboles y profundidad máxima.
  \item Problema: $\min_{j,\lambda\in\Lambda^{(j)}}$ de la pérdida media de validación cruzada.
\end{enumerate}
\resultado \textbf{árbol de profundidad 3 (algoritmo, núcleo, $\gamma$). Activos: $k$ y distancia; núcleo y $C$, más $\gamma$ si el núcleo es gaussiano; número de árboles y profundidad.}
\comprobacion ningún hiperparámetro está activo en dos ramas a la vez.
\end{solucion}

\begin{solucion}{eje:3-sh}
\datos $r_{\min}=2$, $\eta=2$, $r_{\max}=10$; escalones $2,4,8,10$; $N=16$.
\planteamiento \cref{def:sh} con el último escalón recortado.
\pasos
\begin{enumerate}[label=\arabic*.]
  \item (a) 16 con 2 épocas, 8 con 4, 4 con 8 y 2 con 10.
  \item (b) $16\cdot2+8\cdot4+4\cdot8+2\cdot10=32+32+32+20=116$ épocas.
  \item (c) $16\cdot10=160$ épocas: ahorro del 27{,}5\,\%.
  \item (d) $r_{\max}/r_{\min}=5$ no es potencia de $\eta$: el último escalón no multiplica por $\eta$ y no reduce a una sola configuración, y con pocos escalones el cociente $(K+1)/\eta^{K}$ es poco favorable.
\end{enumerate}
\resultado \textbf{116 épocas frente a 160.}
\comprobacion los tres primeros escalones cuestan $Nr_{\min}=32$ cada uno, como en la \cref{prop:coste-sh}.
\end{solucion}

\begin{solucion}{eje:3-celdas}
\datos $B=4$, $\abs{\mathcal{O}}=5$.
\planteamiento fórmula $\prod_{b=1}^{B}((b+1)\abs{\mathcal{O}})^{2}$.
\pasos
\begin{enumerate}[label=\arabic*.]
  \item (a) $(10\cdot15\cdot20\cdot25)^{2}=75\,000^{2}=5{,}625\cdot10^{9}$.
  \item (b) Dos celdas independientes: $(5{,}625\cdot10^{9})^{2}=3{,}16\cdot10^{19}$.
  \item (c) $3{,}16\cdot10^{19}$~h, unos $3{,}6\cdot10^{15}$ años de GPU.
\end{enumerate}
\resultado \textbf{$5{,}6\cdot10^{9}$ celdas; $3{,}2\cdot10^{19}$ redes; enumeración imposible.}
\comprobacion para el nodo $b=1$ hay $2\cdot5=10$ opciones por entrada, y $10^{2}=100$ por nodo.
\end{solucion}

\begin{solucion}{eje:3-darts}
\datos $\alpha=(1{,}0;\ 0{,}2;\ -0{,}5;\ 0{,}8)$; $x=1{,}5$; operaciones $x$, $x^{2}$, $0$, $2x$.
\planteamiento \cref{def:darts}; derivada de la softmax $\partial w_o/\partial\alpha_1=w_o(\Ind[o=1]-w_1)$.
\pasos
\begin{enumerate}[label=\arabic*.]
  \item (a) $e^{\alpha}=(2{,}718;\ 1{,}221;\ 0{,}6065;\ 2{,}226)$, suma $6{,}772$; $w=(0{,}4014;\ 0{,}1804;\ 0{,}0896;\ 0{,}3287)$. $\bar o(1{,}5)=0{,}4014\cdot1{,}5+0{,}1804\cdot2{,}25+0+0{,}3287\cdot3=0{,}6021+0{,}4058+0{,}9860=1{,}9939$.
  \item (b) Mayor $\alpha$: identidad.
  \item (c) $\partial\bar o/\partial\alpha_1=\sum_o o(x)w_o(\Ind[o=1]-w_1)=w_1\bigl(o_1(x)-\bar o(x)\bigr)=0{,}4014\,(1{,}5-1{,}9939)=-0{,}1983$.
\end{enumerate}
\resultado \textbf{$w=(0{,}401;\ 0{,}180;\ 0{,}090;\ 0{,}329)$, $\bar o(1{,}5)=1{,}994$; se conserva la identidad; derivada $-0{,}198$.}
\comprobacion los pesos suman 1; la derivada es negativa porque la identidad da menos que la mezcla en $x=1{,}5$.
\end{solucion}

\begin{solucion}{eje:3-coste}
\datos 12~min proporcionales más 1~min fijo por entrenamiento completo; 0{,}60~€/h.
\planteamiento sumar tiempos y aplicar la \cref{prop:coste-sh} con coste fijo.
\pasos
\begin{enumerate}[label=\arabic*.]
  \item (a) $4^{3}=64$ entrenamientos de 13~min: 832~min $=13{,}87$~h, $8{,}32$~€.
  \item (b) 27 de 13~min: 351~min $=5{,}85$~h, $3{,}51$~€.
  \item (c) $27(1+12/9)+9(1+12/3)+3(1+12)=63+45+39=147$~min $=2{,}45$~h, $1{,}47$~€.
  \item (d) $147/351=0{,}419$: el coste fijo se paga $27+9+3=39$ veces en el halving y 27 en la búsqueda aleatoria.
\end{enumerate}
\resultado \textbf{$8{,}32$; $3{,}51$ y $1{,}47$~€; cociente $0{,}42$ en lugar de $1/3$.}
\comprobacion sin coste fijo: $27\cdot12/9+9\cdot12/3+3\cdot12=36+36+36=108$~min y $108/324=1/3$.
\end{solucion}

\begin{solucion}{eje:3-cruce}
\datos $f_1(r)=0{,}30+0{,}20e^{-r/2}$, $f_2(r)=0{,}25+0{,}40e^{-r/8}$.
\planteamiento resolver $f_1=f_2$; aplicar la \cref{prop:acierto-sh}.
\pasos
\begin{enumerate}[label=\arabic*.]
  \item (a) $g(r)=f_1-f_2=0{,}05+0{,}2e^{-r/2}-0{,}4e^{-r/8}$; $g(16)=0{,}0500+0{,}0001-0{,}0541=-0{,}0041$ y $g(17)=0{,}0500+0{,}0000-0{,}0477=+0{,}0023$; por bisección, $r^{*}\approx16{,}6$ épocas.
  \item (b) Escalón 0 ($r=16$): $f_1=0{,}3001<f_2=0{,}3041$; se conserva $\lambda_1$, que con $r=32$ vale $0{,}3000$, mientras que $f_2(32)=0{,}25+0{,}4e^{-4}=0{,}2573$. Se devuelve la peor.
  \item (c) Usar $r_{\min}\ge17$, o ejecutar como Hyperband una ronda sin eliminación ($r_{\min}=r_{\max}$), que evalúa ambas a 32 épocas: coste $2\cdot32=64$ épocas frente a $2\cdot16+32=64$ del halving; aquí no cuesta más. En general, Hyperband multiplica el coste por el número de rondas, a cambio de no depender de una sola $r_{\min}$.
\end{enumerate}
\resultado \textbf{$r^{*}\approx16{,}6$; el halving con $r_{\min}=16$ elige $\lambda_1$ (pérdida final $0{,}300$ frente a $0{,}257$).}
\comprobacion $g$ es creciente para $r>4$ (su derivada $-0{,}1e^{-r/2}+0{,}05e^{-r/8}$ es positiva cuando $e^{3r/8}>2$), luego el cruce es único.
\end{solucion}

% ---------------------------------------------------------------------
\section{Tema 4}

\begin{solucion}{eje:4-deriva}
\datos cuatro situaciones.
\planteamiento \cref{def:deriva}.
\pasos
\begin{enumerate}[label=\arabic*.]
  \item (a) Cambia $P(X)$ (vibraciones), no $P(Y\mid X)$: cambio de covariables.
  \item (b) Con la misma $X$ cambia la probabilidad de avería: deriva de concepto.
  \item (c) Sube $P(Y)$ con $P(X\mid Y)$ igual: cambio de etiquetas.
  \item (d) Cambia la definición de la característica: cambio de la característica (sesgo del sistema, \cref{def:skew}), que debe corregirse en la ingesta, no reentrenando.
\end{enumerate}
\resultado \textbf{(a) covariables; (b) concepto; (c) etiquetas; (d) cambio de unidad de una característica.}
\comprobacion solo en (b) un modelo perfecto con los datos antiguos dejaría de ser correcto aunque se le dieran entradas bien medidas.
\end{solucion}

\begin{solucion}{eje:4-deuda}
\datos cuatro situaciones.
\planteamiento \cref{def:deuda}.
\pasos
\begin{enumerate}[label=\arabic*.]
  \item (a) Consumidores no declarados.
  \item (b) Entrelazamiento (CACE).
  \item (c) Deuda de configuración.
  \item (d) Cambios en el mundo exterior: el umbral óptimo $c_p/c_f$ depende de costes que cambian.
\end{enumerate}
\resultado \textbf{(a) consumidores no declarados; (b) CACE; (c) configuración; (d) mundo exterior.}
\comprobacion cada caso tiene un remedio distinto: registro de consumidores, pruebas de regresión, revisión de configuración, umbrales calculados a partir de costes actualizados.
\end{solucion}

\begin{solucion}{eje:4-uniforme}
\datos $T\sim U(0,1500)$ días.
\planteamiento \cref{def:fiabilidad}.
\pasos
\begin{enumerate}[label=\arabic*.]
  \item (a) $F(t)=t/1500$, $R(t)=1-t/1500$, $f=1/1500$, $\lambda(t)=1/(1500-t)$ para $0\le t<1500$.
  \item (b) $\lambda(0)=6{,}67\cdot10^{-4}$, $\lambda(500)=1{,}00\cdot10^{-3}$ y $\lambda(1000)=2{,}00\cdot10^{-3}$ por día.
  \item (c) No: la tasa de fallo crece (vida IFR); a los 500 días es un 50\,\% mayor que nueva. Además, la vida restante es uniforme en $(0,1000)$, de media 500 días, frente a 750 de una nueva.
  \item (d) Solo para la exponencial, de tasa constante y sin memoria.
\end{enumerate}
\resultado \textbf{$\lambda(t)=1/(1500-t)$, creciente: la cinta usada no es tan buena como nueva.}
\comprobacion $R(t)=\exp(-\int_0^t\frac{du}{1500-u})=\exp(\ln\frac{1500-t}{1500})=1-t/1500$.
\end{solucion}

\begin{solucion}{eje:4-psi}
\datos $p=(0{,}10;0{,}40;0{,}40;0{,}10)$, $q=(0{,}05;0{,}30;0{,}45;0{,}20)$.
\planteamiento \cref{def:psi} y \cref{prop:psi}.
\pasos
\begin{enumerate}[label=\arabic*.]
  \item (a) Términos $(q_b-p_b)\ln(q_b/p_b)$: $0{,}0347$; $0{,}0288$; $0{,}0059$; $0{,}0693$. $\PSI=0{,}1386$.
  \item (b) $\KL(p\|q)=0{,}1\ln2+0{,}4\ln\frac43+0{,}4\ln\frac89+0{,}1\ln\frac12=0{,}0680$; $\KL(q\|p)=0{,}0707$; suma $0{,}1386$.
  \item (c) Acumuladas: $P=(0{,}10;0{,}50;0{,}90;1)$, $Q=(0{,}05;0{,}35;0{,}80;1)$; diferencias $0{,}05$; $0{,}15$; $0{,}10$; $0$: estadístico $0{,}15$.
\end{enumerate}
\resultado \textbf{$\PSI=0{,}139$; $\KL(p\|q)=0{,}068$ y $\KL(q\|p)=0{,}071$; KS $=0{,}15$.}
\comprobacion $0{,}0680+0{,}0707=0{,}1387\approx0{,}1386$ (redondeo).
\end{solucion}

\begin{solucion}{eje:4-edad}
\datos exponencial de media 200~d; $c_p=500$, $c_f=3000$; Weibull $\beta=2$, $\theta=100$.
\planteamiento \cref{prop:edad}.
\pasos
\begin{enumerate}[label=\arabic*.]
  \item (a) $\int_0^{T}e^{-u/200}du=200(1-e^{-T/200})$. $T=50$: $R=0{,}7788$, $C=(500\cdot0{,}7788+3000\cdot0{,}2212)/44{,}24=23{,}80$~€/d. $T=100$: $R=0{,}6065$, $C=(303{,}3+1180{,}6)/78{,}69=18{,}85$. $T=\infty$: $3000/200=15{,}00$.
  \item (b) Con vida exponencial la máquina usada es como nueva; intervenir antes solo añade costes y acorta los ciclos sin reducir la tasa de fallo: $T^{*}=\infty$, como predice la proposición para $\beta=1$.
  \item (c) $c_p/(c_f-c_p)=0{,}2$; $\lambda(T)=2T/10^{4}$. $G(40)=0{,}008\cdot37{,}97-(1-e^{-0{,}16})=0{,}3037-0{,}1479=0{,}156$; $G(50)=0{,}01\cdot46{,}13-(1-e^{-0{,}25})=0{,}4613-0{,}2212=0{,}240$. Como $G$ es creciente y pasa de $0{,}156$ a $0{,}240$, $T^{*}\in(40,50)$.
\end{enumerate}
\resultado \textbf{(a) $23{,}80$; $18{,}85$; $15{,}00$~€/día: no conviene el preventivo; (c) $40<T^{*}<50$ días (numéricamente $T^{*}=45{,}5$ y $C(T^{*})=2500\lambda(T^{*})=22{,}74$~€/día).}
\comprobacion $C$ decrece en (a) ($23{,}80>18{,}85>15{,}00$), como exige el caso DFR (aquí, tasa constante).
\end{solucion}

\begin{solucion}{eje:4-ips}
\datos $\pi_0(\text{interv.}\mid\text{alarma})=0{,}3$; $\pi_0(\text{interv.}\mid\text{sin})=0{,}1$; seis registros.
\planteamiento \cref{prop:ips} con $\pi$ determinista.
\pasos
\begin{enumerate}[label=\arabic*.]
  \item Pesos $\pi/\pi_0$: registro 1 (alarma, interviene): $1/0{,}3=3{,}333$; 2 y 3 (alarma, no): $0$; 4 y 6 (sin alarma, no): $1/0{,}9=1{,}111$; 5 (sin alarma, interviene): $0$.
  \item Suma de pesos por recompensa: $3{,}333\cdot(-1{,}2)+1{,}111\cdot0+1{,}111\cdot0=-4{,}0$.
  \item $\hat V_{\mathrm{IPS}}(\pi)=-4{,}0/6=-0{,}667$ miles de euros por semana.
  \item Política antigua (media de los registros): $(-1{,}2-9+0+0-1{,}2+0)/6=-1{,}9$.
\end{enumerate}
\resultado \textbf{$\hat V(\pi)\approx-0{,}67$ frente a $-1{,}9$ (miles de euros); la estimación favorece la política nueva, pero descansa en un único registro con alarma e intervención, multiplicado por $3{,}33$: es muy poco fiable.}
\comprobacion los pesos de los registros coherentes con $\pi$ suman $3{,}333+1{,}111+1{,}111=5{,}56$, cerca de $n=6$, como se espera de media.
\end{solucion}

\begin{solucion}{eje:4-reentreno}
\datos $c_r=3000$~€; $a=4$~€/día$^{2}$.
\planteamiento \cref{prop:reentreno}.
\pasos
\begin{enumerate}[label=\arabic*.]
  \item (a) $T^{*}=\sqrt{6000/4}=38{,}73$ días; coste $\sqrt{2\cdot4\cdot3000}=154{,}9$~€/día.
  \item (b) $T=7$: $3000/7+4\cdot7/2=428{,}6+14=442{,}6$~€/día.
  \item (c) $T=180$: $16{,}7+360=376{,}7$~€/día.
\end{enumerate}
\resultado \textbf{$T^{*}\approx39$ días con $154{,}9$~€/día; semanal $442{,}6$ y semestral $376{,}7$~€/día.}
\comprobacion en $T^{*}$ los dos términos coinciden: $3000/38{,}73=77{,}46=4\cdot38{,}73/2$.
\end{solucion}

\begin{solucion}{eje:4-lineabase}
\datos tasa de positivos $\pi=0{,}125$.
\planteamiento calcular precisión, sensibilidad, $F_1$ y exactitud de predictores que ignoran $X$.
\pasos
\begin{enumerate}[label=\arabic*.]
  \item (a) Predice positivo con probabilidad $0{,}125$: precisión $=\pi=0{,}125$, sensibilidad $=0{,}125$, $F_1=0{,}125$; exactitud $\pi^{2}+(1-\pi)^{2}=0{,}0156+0{,}7656=0{,}781$.
  \item (b) Con probabilidad $1/2$: precisión $0{,}125$, sensibilidad $0{,}5$, $F_1=2\cdot0{,}125\cdot0{,}5/0{,}625=0{,}2$; exactitud $0{,}5$.
  \item (c) Siempre negativo: exactitud $0{,}875$; sensibilidad $0$ y $F_1=0$.
\end{enumerate}
\resultado \textbf{(a) exactitud $0{,}781$, $F_1=0{,}125$; (b) $0{,}5$ y $0{,}2$; (c) $0{,}875$ y $0$.}
\comprobacion la mejor exactitud la da un predictor inútil; por eso se informa de $F_1$, del área ROC, de la pérdida logarítmica y de la calibración.
\end{solucion}

\begin{solucion}{eje:4-monitor}
\datos optimizador del \cref{ej:caso-mant}; etiquetas con 7 días de retraso.
\planteamiento combinar la \cref{def:deriva}, la \cref{prop:psi}, la \cref{prop:reentreno} y la \cref{obs:bucle}.
\pasos
\begin{enumerate}[label=\arabic*.]
  \item Diario: esquema y ausencias de los sensores; detección de sensores congelados (varianza nula en 24~h); distribución de las probabilidades predichas; latencia y errores del servicio. Acción: alarma al responsable de datos y paso a la política por edad para la máquina afectada.
  \item Semanal: PSI por variable y por máquina frente a una referencia de 8 semanas estables, con umbral calibrado como el percentil 99 del PSI entre semanas estables; tasa de anulaciones humanas y sus motivos; plan frente a ejecución.
  \item Mensual (con etiquetas): pérdida logarítmica, calibración por quintiles y área ROC de las predicciones de hace al menos 7 días; averías y coste frente al gemelo; decisión de reentrenar por periodo ($T^{*}$) o por disparador.
  \item Bucle: cada decisión se registra con $p_7$, la acción y quién la aprobó; las intervenciones se tratan como censura; el reentrenamiento usa solo datos de máquinas sin intervención en la ventana de la etiqueta, o modelos de supervivencia con censura.
\end{enumerate}
\resultado \textbf{plan en tres frecuencias con estadísticos, umbrales calibrados, acciones definidas y tratamiento de la censura.}
\comprobacion cubre entradas, predicciones, decisiones, resultados y salud del servicio, como pide la clase~10 de CS329S~\cite{cs329s}.
\end{solucion}

% ---------------------------------------------------------------------
\section{Tema 5}

\begin{solucion}{eje:5-alternativa}
\datos ahorro declarado del 18\,\% frente a las rutas de los conductores.
\planteamiento la alternativa honrada es la mejor solución sin IA (\cref{obs:errores5}).
\pasos
\begin{enumerate}[label=\arabic*.]
  \item Alternativas: un resolvedor clásico de rutas con ventanas horarias sin predicción (módulo~24, tema~3); la misma optimización con tiempos medios históricos por franja; reglas operativas mejoradas (agrupar por zonas, secuencias fijas).
  \item Parte no atribuible a la IA: la que se obtiene solo por optimizar las rutas con un resolvedor (pasar de rutas manuales a rutas optimizadas) y la debida a cambios simultáneos (menos pedidos, otra flota).
\end{enumerate}
\resultado \textbf{comparar con un resolvedor sin predicción y con reglas mejoradas; la IA solo aporta la diferencia con la mejor de ellas.}
\comprobacion en el caso de mantenimiento, la IA aporta $34\,370$~€/año de los $170\,990$ de mejora total.
\end{solucion}

\begin{solucion}{eje:5-pareado}
\datos $\sigma_B=6500$, $\sigma_D=5200$; $r=0{,}78$.
\planteamiento \cref{prop:pareado}.
\pasos
\begin{enumerate}[label=\arabic*.]
  \item (a) Independientes: $\sqrt{6500^{2}+5200^{2}}=\sqrt{6{,}929\cdot10^{7}}=8324$~€. Pareadas:
  \[
  \sqrt{6{,}929\cdot10^{7}-2\cdot0{,}78\cdot6500\cdot5200}=\sqrt{1{,}656\cdot10^{7}}=4070\ \text{€}.
  \]
  \item (b) Para la misma precisión de la media, $n$ es proporcional a la varianza: $6{,}929/1{,}656=4{,}18$ veces más semanas.
\end{enumerate}
\resultado \textbf{8324~€ frente a 4070~€; harían falta unas 4{,}2 veces más semanas sin emparejar.}
\comprobacion $8324^{2}/4070^{2}=4{,}18$.
\end{solucion}

\begin{solucion}{eje:5-bootstrap}
\datos $n=10$, $\bar d=40{,}3$; segunda y decimoctava medias \emph{bootstrap}: $39{,}4$ y $42{,}4$; $s=4{,}67$; $t_{9;0{,}90}=1{,}383$.
\planteamiento \cref{def:bootstrap} con $\alpha=0{,}2$ y 20 remuestras (cuantiles $0{,}1$ y $0{,}9$ = valores 2.º y 18.º).
\pasos
\begin{enumerate}[label=\arabic*.]
  \item (a) Percentil: $[39{,}4;\ 42{,}4]$.
  \item (b) $\delta^{*}_{(0{,}1)}=39{,}4-40{,}3=-0{,}9$ y $\delta^{*}_{(0{,}9)}=42{,}4-40{,}3=2{,}1$; básico: $[40{,}3-2{,}1;\ 40{,}3+0{,}9]=[38{,}2;\ 41{,}2]$.
  \item (c) $40{,}3\pm1{,}383\cdot4{,}67/\sqrt{10}=40{,}3\pm2{,}04$: $[38{,}26;\ 42{,}34]$.
  \item (d) Con 20 remuestras los cuantiles dependen de dos valores y varían mucho entre ejecuciones; con 10\,000, el error de Monte Carlo del cuantil es despreciable.
\end{enumerate}
\resultado \textbf{percentil $[39{,}4;42{,}4]$; básico $[38{,}2;41{,}2]$; $t$: $[38{,}3;42{,}3]$ (en cientos de euros por semana).}
\comprobacion el intervalo básico es el percentil reflejado respecto a $\bar d$: $40{,}3\pm(42{,}4-40{,}3)$ y $40{,}3\pm(40{,}3-39{,}4)$ intercambiados.
\end{solucion}

\begin{solucion}{eje:5-van}
\datos $I=48\,000$~€; flujo neto $25\,370$~€/año; $N=5$; $\rho=8\,\%$.
\planteamiento \cref{prop:anualidad}.
\pasos
\begin{enumerate}[label=\arabic*.]
  \item (a) $(1-1{,}08^{-5})/0{,}08=(1-0{,}68058)/0{,}08=3{,}9927$.
  \item (b) $a_{5;0{,}12}=(1-1{,}12^{-5})/0{,}12=3{,}6048$; $\VAN=25\,370\cdot3{,}6048-48\,000=43\,453$~€.
  \item (c) Flujo neto mínimo $48\,000/3{,}9927=12\,022$~€, es decir, beneficio bruto de $12\,022+9000=21\,022$~€/año.
  \item (d) $a_{3;0{,}08}=2{,}5771$; $\VAN=25\,370\cdot2{,}5771-48\,000=17\,381$~€.
\end{enumerate}
\resultado \textbf{(a) $3{,}9927$; (b) $43\,453$~€; (c) $21\,022$~€/año de beneficio; (d) $17\,381$~€.}
\comprobacion con $\rho=8\,\%$ y $N=5$: $25\,370\cdot3{,}9927-48\,000=53\,295$, el valor del \cref{ej:caso-valor}.
\end{solucion}

\begin{solucion}{eje:5-tamano}
\datos $\sigma_D=4356$, $\delta=661$; parada no planificada: $\sigma_D=3{,}1$~h, $\delta=0{,}9$~h.
\planteamiento \cref{prop:tamano}.
\pasos
\begin{enumerate}[label=\arabic*.]
  \item (a) $z_{0{,}975}+z_{0{,}9}=1{,}960+1{,}282=3{,}242$; $n=(3{,}242\cdot4356/661)^{2}=21{,}36^{2}=456{,}3$: 457 semanas.
  \item (b) Con potencia $0{,}8$ hacen falta 341 pares; con 6 líneas en paralelo, $341/6=56{,}8$: 57 semanas.
  \item (c) $(2{,}802\cdot3{,}1/0{,}9)^{2}=9{,}65^{2}=93{,}1$: 94 semanas.
\end{enumerate}
\resultado \textbf{(a) 457; (b) 57; (c) 94 semanas.}
\comprobacion la relación señal-ruido $\delta/\sigma_D$ pasa de $0{,}15$ (coste) a $0{,}29$ (horas de parada), y el tamaño se divide por $(0{,}29/0{,}15)^{2}\approx3{,}7$.
\end{solucion}

\begin{solucion}{eje:5-cvar}
\datos costes ordenados $139,148,149,152,157,160,166,171,175,183$ (miles de euros).
\planteamiento VaR y CVaR empíricos; \cref{prop:cvar} con la distribución empírica.
\pasos
\begin{enumerate}[label=\arabic*.]
  \item (a) Media $1600/10=160$. VaR$_{0{,}8}$: el menor valor con frecuencia acumulada $\ge0{,}8$, el octavo, $171$. CVaR$_{0{,}8}$: media de los dos peores, $(175+183)/2=179$.
  \item (b) $\phi(\tau)=\tau+\frac{1}{0{,}2}\cdot\frac{1}{10}\sum_s(c_s-\tau)^{+}$. $\phi(175)=175+0{,}5\cdot8=179$. $\phi(160)=160+0{,}5\cdot(6+11+15+23)=187{,}5$. $\phi(183)=183+0=183$.
\end{enumerate}
\resultado \textbf{media 160; VaR 171; CVaR 179; $\phi(175)=179<\phi(183)=183<\phi(160)=187{,}5$.}
\comprobacion con distribución discreta, $\phi$ es constante e igual a 179 en todo $[171,175]$; el mínimo de la \cref{prop:cvar} no es único, pero su valor sí.
\end{solucion}

\begin{solucion}{eje:5-amfe}
\datos cinco modos del \cref{tab:amfe} y dos nuevos.
\planteamiento \cref{def:amfe}.
\pasos
\begin{enumerate}[label=\arabic*.]
  \item NPR nuevos: ciberataque $9\cdot2\cdot8=144$; datos de vacaciones $5\cdot3\cdot5=75$.
  \item Orden: sensor congelado (168) y confianza excesiva (168); etiquetas mal registradas (144) y ciberataque (144, con gravedad 9); deriva (120); vacaciones (75); resolvedor (20).
  \item Medidas para los primeros: sensor congelado, regla de varianza nula y redundancia del sensor de vibración en la llenadora; confianza excesiva, mostrar motivos y probabilidades, auditoría mensual de órdenes aceptadas sin cambios y formación sobre el sesgo de automatización.
\end{enumerate}
\resultado \textbf{NPR 144 y 75; prioridad: sensor congelado y confianza excesiva; el ciberataque se trata en todo caso por su gravedad.}
\comprobacion la gravedad 9 del ciberataque obliga a tratarlo aunque su NPR empate con otros (\cref{def:amfe}).
\end{solucion}

\begin{solucion}{eje:5-reglamento}
\datos tres sistemas.
\planteamiento artículo~6 y anexo~III del Reglamento~\cite{aiact}, en términos generales; no es una calificación jurídica.
\pasos
\begin{enumerate}[label=\arabic*.]
  \item (a) Optimizador de mantenimiento de una línea de envasado: no gestiona una infraestructura crítica ni actúa sobre funciones de seguridad de las máquinas; un técnico revisa sus órdenes. No parece de alto riesgo; faltaría confirmar que no se integra en un producto regulado como componente de seguridad. Medidas recomendables: registro, supervisión humana, documentación y monitorización.
  \item (b) Válvulas de una red de calefacción urbana para evitar sobrepresiones: puede ser un componente de seguridad en la operación del suministro de calefacción (anexo~III, punto~2), por tanto candidato a alto riesgo. Faltaría saber si es la función de seguridad o una capa de optimización sobre un sistema de seguridad independiente (módulo~26, tema~6). Medidas: todas las de los artículos~9 a~15, con especial atención a la robustez y la supervisión del artículo~14.
  \item (c) Asistente que redacta borradores revisados: tarea preparatoria sin influencia sustancial en decisiones de alto riesgo; en general no es de alto riesgo. Medidas: evaluación de calidad, registro, revisión humana obligatoria antes de emitir la orden.
\end{enumerate}
\resultado \textbf{(a) y (c) previsiblemente no; (b) posible alto riesgo según su función; en todos, adoptar registro, supervisión y documentación.}
\comprobacion el criterio decisivo es la función de seguridad y el ámbito del anexo~III, no la técnica de IA empleada.
\end{solucion}

% ---------------------------------------------------------------------
\section{Tema 6}

\begin{solucion}{eje:6-orden}
\datos propuesta «modelo primero».
\planteamiento comparar con los doce pasos.
\pasos
\begin{enumerate}[label=\arabic*.]
  \item Faltan: la decisión y sus responsables (1), el objetivo y las restricciones (2), la línea base y el techo (3), el inventario de datos (4) y la formulación (5); la medida del ahorro se deja al final sin diseño (9).
  \item Orden correcto de las cinco primeras tareas: definir la decisión y quién la toma; fijar objetivo, indicadores y restricciones; medir la línea base, la mejor alternativa sencilla y el techo; inventariar los datos y su disponibilidad; formular el problema y elegir la familia de métodos. Solo entonces se elige el modelo predictivo (que puede no ser un transformer).
\end{enumerate}
\resultado \textbf{la propuesta salta los pasos 1--5 y deja la evaluación sin diseñar.}
\comprobacion la arquitectura (nube) y el modelo (transformer) son decisiones de los pasos 6 y 10.
\end{solucion}

\begin{solucion}{eje:6-margenes}
\datos 2~s; margen del 25\,\%; percentiles 99: $0{,}30$; $0{,}15$; $0{,}05$; $0{,}70$; $0{,}20$~s.
\planteamiento asignación de presupuesto con margen de sistema.
\pasos
\begin{enumerate}[label=\arabic*.]
  \item (a) Presupuesto asignable: $0{,}75\cdot2=1{,}5$~s. Se reparte en proporción a la variabilidad esperada, dando más al resolvedor (tiempo variable) y a la lectura de datos (depende de la red): por ejemplo $0{,}35$; $0{,}15$; $0{,}10$; $0{,}65$; $0{,}25$~s.
  \item (b) Suma de percentiles 99: $1{,}40$~s $\le1{,}5$~s: se cumple con margen ($0{,}60$~s, 30\,\%). La suma de percentiles es conservadora (el percentil 99 de la suma suele ser menor). Revisar el resolvedor: es el mayor y el más variable; fijar un límite de tiempo y arrancar desde el plan anterior.
\end{enumerate}
\resultado \textbf{1{,}5~s asignables; las mediciones suman 1{,}40~s y cumplen; vigilar el resolvedor.}
\comprobacion $0{,}35+0{,}15+0{,}10+0{,}65+0{,}25=1{,}50$.
\end{solucion}

\begin{solucion}{eje:6-alternativas}
\datos pesos $(0{,}4;0{,}2;0{,}2;0{,}2)$; I $(2,5,5,4)$, II $(4,4,3,4)$, III $(5,2,1,2)$.
\planteamiento suma ponderada y dominancia de Pareto (módulo~24, tema~5).
\pasos
\begin{enumerate}[label=\arabic*.]
  \item (a) I: $0{,}8+1+1+0{,}8=3{,}6$; II: $1{,}6+0{,}8+0{,}6+0{,}8=3{,}8$; III: $2{,}0+0{,}4+0{,}2+0{,}4=3{,}0$.
  \item (b) Ninguno domina a otro: III es el mejor en coste; I, en interpretabilidad y esfuerzo; II es mejor que I en coste y que III en todo lo demás. Los tres son de Pareto.
  \item (c) Con peso $w$ para el coste y $(1-w)/3$ para cada uno de los otros: II $=4w+\frac{11}{3}(1-w)$, III $=5w+\frac53(1-w)$. Igualando: $4w+\frac{11}3-\frac{11}3w=5w+\frac53-\frac53w$, es decir, $-3w=-2$ y $w=2/3$.
\end{enumerate}
\resultado \textbf{II gana con $3{,}8$; los tres son de Pareto; III empataría con II si el coste pesara $2/3$.}
\comprobacion con $w=2/3$: II $=8/3+11/9=35/9$ y III $=10/3+5/9=35/9$.
\end{solucion}

\begin{solucion}{eje:6-techo}
\datos $330\,290$; $193\,670$; $159\,300$; $76\,170$~€/año; sensor: 30\,000~€ y 3000~€/año; 5 años al 8\,\%.
\planteamiento fracciones del margen; \cref{prop:anualidad}.
\pasos
\begin{enumerate}[label=\arabic*.]
  \item (a) Margen total $330\,290-76\,170=254\,120$. Política por edad: $136\,620/254\,120=53{,}8\,\%$. Optimizador: $170\,990/254\,120=67{,}3\,\%$.
  \item (b) Margen restante $159\,300-76\,170=83\,130$. VAN $\ge0\iff(b-3000)\cdot3{,}9927\ge30\,000\iff b\ge10\,514$~€/año, un $12{,}6\,\%$ del margen restante.
\end{enumerate}
\resultado \textbf{53{,}8\,\% y 67{,}3\,\%; el sensor debe capturar al menos un 12{,}6\,\% del margen restante.}
\comprobacion $(10\,514-3000)\cdot3{,}9927=30\,001$.
\end{solucion}

\begin{solucion}{eje:6-formatos}
\datos matriz de cambios del \cref{ej:mini-fabrica}.
\planteamiento enumerar los seis ciclos.
\pasos
\begin{enumerate}[label=\arabic*.]
  \item (a) Prohibir D$\to$C elimina los ciclos A--B--D--C (215) y A--D--C--B (210). El óptimo A--B--C--D (D$\to$A) sigue permitido: no cambia.
  \item (b) Cada ciclo contiene exactamente un cambio desde D; con 60~min se restan 30 a todos: 145, 160, 180, 185, 195 y 205 para A--B--C--D, A--C--D--B, A--D--C--B, A--B--D--C, A--D--B--C y A--C--B--D. El orden no cambia.
  \item (c) Con 8 formatos y el primero fijo, $7!=5040$ ciclos.
\end{enumerate}
\resultado \textbf{(a) no cambia (175~min); (b) 145~min con el mismo orden; (c) 5040 ciclos.}
\comprobacion $175-30=145$.
\end{solucion}

\begin{solucion}{eje:6-lod}
\datos tres objetos con tres LOD; presupuesto 100 (miles de triángulos).
\planteamiento enumeración de las 27 combinaciones; mochila de elección múltiple.
\pasos
\begin{enumerate}[label=\arabic*.]
  \item (a) Factibles con menor error sin regla: (LOD1, LOD1, LOD1) con 70 y error $2{,}4$; las de error menor, como (0,1,1) con error $1{,}4$ y (1,1,0) con $1{,}6$, usan 105 y no caben. Con el personaje en LOD0: (0,2,1) con 87 y error $3{,}3$; (0,1,2) con 95 y $3{,}6$; (0,2,2) con 77 y $5{,}5$.
  \item (b) La regla no cuesta nada cuando el óptimo sin regla ya tiene el personaje en LOD0: con presupuesto 105 el óptimo es (0,1,1), error $1{,}4$. Por debajo de 105, el óptimo sin regla es (1,1,1) o (1,1,0), con el personaje en LOD1. Presupuesto mínimo: 105\,000 triángulos.
  \item (c) $\min\sum_{o,k}e_{ok}z_{ok}$ s.a. $\sum_kz_{ok}=1$ para cada objeto, $\sum_{o,k}t_{ok}z_{ok}\le B$, $z_{\text{personaje},k}=0$ para $k\ge1$, $z_{ok}\in\{0,1\}$.
\end{enumerate}
\resultado \textbf{sin regla: error $2{,}4$; con regla: $3{,}3$; la regla es gratuita desde 105\,000 triángulos.}
\comprobacion (0,2,1): $60+12+15=87\le100$ y $0+2{,}5+0{,}8=3{,}3$.
\end{solucion}

\begin{solucion}{eje:6-simpson}
\datos planta 1: $81/87$ y $234/270$; planta 2: $192/263$ y $55/80$.
\planteamiento tasas por segmento y globales.
\pasos
\begin{enumerate}[label=\arabic*.]
  \item (a) Planta 1: optimizador $93{,}1\,\%$, anterior $86{,}7\,\%$. Planta 2: $73{,}0\,\%$ y $68{,}8\,\%$. Global: optimizador $273/350=78{,}0\,\%$, anterior $289/350=82{,}6\,\%$.
  \item (b) El optimizador se evaluó sobre todo en la planta 2 (263 de 350 casos), donde las averías son más difíciles de evitar; la política anterior, sobre todo en la planta 1. La mezcla, no la calidad, invierte el resultado global: paradoja de Simpson.
  \item (c) Presentar los resultados por segmento, con el tamaño de cada uno, y comparar con la misma mezcla (o con diseño pareado por planta).
\end{enumerate}
\resultado \textbf{el optimizador gana en cada planta y pierde en el total por la mezcla de casos.}
\comprobacion $81+192=273$ y $234+55=289$, ambos sobre 350.
\end{solucion}

\begin{solucion}{eje:6-diseno}
\datos almacén frigorífico de 20\,000~m$^{3}$; tarifa horaria; placas solares; $-22\le T\le-18$~ºC.
\planteamiento recorrer el \cref{tab:plantilla}.
\pasos
\begin{enumerate}[label=\arabic*.]
  \item Decisión (1--2): consignas horarias de los compresores para las próximas 24~h; responsable: jefe de mantenimiento; puede pararlo el operador de turno.
  \item Objetivo y requisitos (3): minimizar el coste eléctrico; restricción dura $-22\le T\le-18$~ºC en todas las sondas de producto; indicadores: kWh/día, horas cerca del límite, arranques de compresor. Módulo~24, tema~1.
  \item Línea base y techo (4): coste actual con termostato fijo; techo con información perfecta de precios y radiación. Primer número: coste diario actual.
  \item Datos (5): sondas, contadores, precios, predicción de radiación; versionado. Módulo~28, tema~4.
  \item Formulación (6): control predictivo con modelo térmico identificado (módulo~26, temas~2--3); programa lineal o cuadrático por hora.
  \item Modelos (7): predicción de la generación solar y de las aperturas de puertas; validación temporal (módulo~8). Primer número: error de predicción de la generación a 24~h.
  \item Salvaguardas (8): termostato existente como respaldo; filtro de seguridad de consignas (módulo~26, tema~6, §6.5).
  \item Evaluación (9--10): gemelo térmico con números aleatorios comunes; piloto en sombra y después días alternos con tamaño calculado (\cref{prop:tamano}).
  \item Despliegue, monitorización, valor y ciclo de vida (11--14): API e interfaz del operador; vigilancia de residuos del modelo térmico; VAN e intervalos (módulo~28, tema~5); revisión anual.
\end{enumerate}
\resultado \textbf{documento completo con un responsable, una restricción dura de temperatura, un MPC lineal, respaldo por termostato y un primer número por sección.}
\comprobacion cada sección de la plantilla queda respondida y enlazada con su módulo.
\end{solucion}
